% =====================================================================
%  CARNET D'ENTRAINEMENT — FICHE 6 — THÈME 5
%  Niveau MOYEN — Corrigé
% =====================================================================
\documentclass[11pt,a4paper]{article}
\usepackage[utf8]{inputenc}
\usepackage[T1]{fontenc}
\usepackage{lmodern}
\usepackage[french]{babel}
\usepackage{amsmath,amssymb,mathtools}
\usepackage{icomma}
\usepackage[version=4]{mhchem}
\usepackage{siunitx}
\sisetup{output-decimal-marker={,}}
\usepackage{xcolor}
\usepackage{colortbl}
\usepackage{tikz}
\usepackage[most]{tcolorbox}
\usepackage{enumitem}
\usepackage{array}
\usepackage{booktabs}
\usepackage{fancyhdr}
\usepackage[hidelinks]{hyperref}
\usetikzlibrary{arrows.meta,positioning,calc,shapes.geometric}
\usepackage[a4paper,margin=2.1cm,headheight=26pt,headsep=8pt]{geometry}

\definecolor{primary}{RGB}{146,95,160}
\definecolor{primarydark}{RGB}{92,52,104}
\definecolor{corcol}{RGB}{0,114,178}
\definecolor{astucecol}{RGB}{198,93,9}
\definecolor{bondcol}{RGB}{60,60,70}
\definecolor{piecol}{RGB}{198,93,9}

\tcbset{boxbase/.style={enhanced,breakable,boxrule=0.9pt,arc=2.5pt,
   left=3mm,right=3mm,top=2.2mm,bottom=2.2mm,
   fonttitle=\bfseries,coltitle=white,toptitle=0.6mm,bottomtitle=0.6mm}}
\newtcolorbox[auto counter]{cor}[1][]{boxbase,colback=corcol!5,colframe=corcol,
   title={\textsc{Corrigé}~\thetcbcounter\ifx\relax#1\relax\relax\else\textmd{ —\itshape\ #1}\fi}}
\newtcolorbox{astucebox}[1][]{boxbase,colback=astucecol!8,colframe=astucecol,
   title={\textsc{Astuce}\ifx\relax#1\relax\relax\else\textmd{ : \itshape #1}\fi}}

\pagestyle{fancy}\fancyhf{}
\fancyhead[L]{\small\textcolor{primarydark}{\textbf{Fiche 6 · Thème 5 — La résonance}}}
\fancyhead[R]{\small\textcolor{primarydark}{\textsc{Moyen · Corrigé}}}
\fancyfoot[C]{\small\thepage}
\renewcommand{\headrulewidth}{0.8pt}
\renewcommand{\headrule}{\hbox to\headwidth{\color{primary}\leaders\hrule height \headrulewidth\hfill}}
\setlength{\parindent}{0pt}\setlength{\parskip}{3pt}

\newcommand{\rep}{\textbf{\textcolor{corcol}{Réponse : }}}

\begin{document}
\begin{center}
{\Large\bfseries\color{primarydark}Thème 5 — Corrigé du niveau Moyen}
\end{center}
\medskip

\begin{cor}[l'ozone \ce{O3}]
\begin{enumerate}[label=\alph*),leftmargin=2em,itemsep=2pt]
  \item \rep les deux formes limites s'écrivent \quad\ce{O=O-O <-> O-O=O}. \\
        On passe de l'une à l'autre en faisant « glisser » la double liaison de l'oxygène de gauche
        vers celui de droite (un doublet libre devient double liaison d'un côté, la double liaison
        redevient doublet libre de l'autre).
  \item \rep les deux formes limites sont \emph{équivalentes} : la double liaison n'appartient
        réellement ni à l'une ni à l'autre position. Les deux liaisons \ce{O-O} sont donc
        \textbf{identiques}. Leur longueur est \textbf{intermédiaire} entre celle d'une liaison
        simple \ce{O-O} (la plus longue) et celle d'une liaison double \ce{O=O} (la plus courte) :
        c'est l'hybride des deux formes.
\end{enumerate}
\end{cor}

\begin{cor}[l'ion carbonate \ce{CO3^2-}]
\begin{enumerate}[label=\alph*),leftmargin=2em,itemsep=2pt]
  \item \rep la double liaison \ce{C=O} peut se placer indifféremment sur \emph{chacun} des trois
        oxygènes. Ces trois positions étant équivalentes, on obtient \textbf{trois} structures
        limites de même énergie, reliées par $\longleftrightarrow$.
  \item \rep les trois liaisons \ce{C-O} sont en réalité \textbf{identiques} (même longueur,
        intermédiaire entre simple et double). Chacune a un « caractère double » partiel de
        $\tfrac{1}{3}$.
  \item \rep par symétrie, la charge $-2$ se répartit \emph{également} sur les trois oxygènes, soit
        \[ \frac{-2}{3} \approx -0{,}67 \text{ par oxygène.} \]
\end{enumerate}
\end{cor}

\begin{cor}[l'ion nitrate \ce{NO3-}]
\begin{enumerate}[label=\alph*),leftmargin=2em,itemsep=2pt]
  \item \rep exactement comme le carbonate, la double liaison \ce{N=O} peut occuper les trois
        positions équivalentes : il y a donc \textbf{3} formes limites équivalentes.
  \item \rep les trois liaisons \ce{N-O} sont \textbf{toutes identiques}, de longueur intermédiaire
        entre une simple \ce{N-O} et une double \ce{N=O}. (La charge $-1$ se répartit sur les trois
        oxygènes, soit $-\tfrac{1}{3}$ chacun.)
\end{enumerate}
\end{cor}

\begin{cor}[deux flèches à ne pas confondre]
\begin{enumerate}[label=\alph*),leftmargin=2em,itemsep=2pt]
  \item \rep \ce{O=O-O <-> O-O=O} met en jeu \textbf{une seule} espèce (l'ozone, décrit par deux
        formules). \ce{2 H2O <=> H3O+ + OH-} met en jeu \textbf{plusieurs} espèces distinctes
        (\ce{H2O}, \ce{H3O+}, \ce{OH-}).
  \item \rep $\longleftrightarrow$ (résonance) signifie « \emph{une même} espèce, dont les électrons
        sont délocalisés, que l'on ne sait pas dessiner d'une seule formule ». $\rightleftharpoons$
        (équilibre) signifie « \emph{deux} espèces réelles qui se transforment l'une en l'autre, et
        coexistent ». Dans le premier cas rien ne se transforme ; dans le second, il y a une vraie
        réaction chimique.
\end{enumerate}
\end{cor}

\begin{cor}[l'ion perchlorate \ce{ClO4-}]
\begin{enumerate}[label=\alph*),leftmargin=2em,itemsep=2pt]
  \item \rep cette structure introduit une \emph{différence} entre les oxygènes (un seul porterait la
        charge $-1$ et une seule liaison simple, les trois autres des doubles liaisons). Or
        l'expérience montre que les quatre oxygènes sont \textbf{équivalents} et les quatre liaisons
        \ce{Cl-O} \textbf{identiques} : la formule unique ne peut donc pas être « la » description
        correcte.
  \item \rep la résonance résout le paradoxe : la charge $-1$ et le « caractère double » ne sont
        localisés sur aucun oxygène en particulier, mais \textbf{délocalisés également} sur les
        quatre. La position de la charge $-1$ « tourne » sur les quatre oxygènes, ce qui donne
        \textbf{4} formes limites équivalentes ; l'hybride est parfaitement symétrique (charge
        moyenne $-\tfrac{1}{4}$ par oxygène).
\end{enumerate}
\end{cor}

\begin{astucebox}[le raccourci « positions équivalentes »]
Pour ces ions à un atome central entouré d'oxygènes équivalents, le nombre de formes limites égale
tout simplement le \textbf{nombre d'oxygènes} pouvant porter la double liaison : 2 pour \ce{O3},
3 pour \ce{CO3^2-} et \ce{NO3-}, 4 pour \ce{ClO4-}. Rappel : seuls les \emph{électrons} bougent,
et $\longleftrightarrow \neq \rightleftharpoons$.
\end{astucebox}

\end{document}
